\documentclass[12pt, a4paper]{article}

\usepackage[utf8]{inputenc}
\usepackage[T1]{fontenc}
\usepackage[spanish]{babel}
\usepackage{geometry}
\usepackage{listings}
\usepackage{xcolor}
\usepackage{titlesec}
\usepackage{booktabs}
\usepackage{tabularx}
\usepackage{hyperref}
\usepackage{fancyhdr}
\usepackage{graphicx}
\usepackage{enumitem}
\usepackage{parskip}

\geometry{margin=2.5cm, headheight=15pt}

% Colores para codigo SQL
\definecolor{sqlblue}{RGB}{0, 0, 180}
\definecolor{sqlgreen}{RGB}{0, 128, 0}
\definecolor{sqlgray}{RGB}{128, 128, 128}
\definecolor{sqlbg}{RGB}{245, 245, 245}
\definecolor{sqlred}{RGB}{180, 0, 0}

% Estilo para bloques de codigo SQL
\lstdefinestyle{sql}{
    language=SQL,
    backgroundcolor=\color{sqlbg},
    basicstyle=\ttfamily\small,
    keywordstyle=\color{sqlblue}\bfseries,
    commentstyle=\color{sqlgreen}\itshape,
    stringstyle=\color{sqlred},
    numberstyle=\tiny\color{sqlgray},
    numbers=left,
    stepnumber=1,
    breaklines=true,
    breakatwhitespace=false,
    frame=single,
    rulecolor=\color{sqlgray},
    tabsize=4,
    showstringspaces=false,
    captionpos=b,
    morekeywords={SEQUENCE, NOCYCLE, NOCACHE, MINVALUE, MAXVALUE, INCREMENT,
                  TRIGGER, BEFORE, AFTER, EACH, ROW, RETURNS, LANGUAGE, plpgsql,
                  DECLARE, BEGIN, END, NEW, OLD, FOUND, RAISE, EXCEPTION, NOTICE,
                  COALESCE, BOOLEAN, NUMERIC, BYTEA, TEXT, CURRENT_DATE,
                  CURRENT_TIME, NEXTVAL, LPAD, SQLERRM, DO, CASCADE, RESTRICT,
                  SET, NULL, REPLACE, FUNCTION, EXECUTE, LATERAL, IF, THEN,
                  ELSIF, ELSE, RETURN, VARCHAR, INTEGER, DATE, TIME}
}

\lstset{style=sql}

% Encabezado y pie de pagina
\pagestyle{fancy}
\fancyhf{}
\rhead{Sistema de Restaurante}
\lhead{Proyecto Final BD --- Grupo 01}
\rfoot{P\'agina \thepage}
\lfoot{PostgreSQL}

% Formato de titulos de seccion
\titleformat{\section}{\large\bfseries}{}{0em}{}[\titlerule]
\titleformat{\subsection}{\normalsize\bfseries}{}{0em}{}

\hypersetup{
    colorlinks=true,
    linkcolor=sqlblue,
    urlcolor=sqlblue
}

% ============================================================
\begin{document}

\begin{titlepage}
    \centering
    \vspace*{2cm}
    {\LARGE \textbf{Sistema de Base de Datos para Restaurante}}\\[0.5cm]
    {\large Documentaci\'on T\'ecnica Completa}\\[2cm]
    {\large Proyecto Final --- Bases de Datos}\\[0.3cm]
    {\large Grupo 01}\\[3cm]
    \begin{tabular}{ll}
        \textbf{Motor:}   & PostgreSQL 18          \\[0.3cm]
        \textbf{Fecha:}   & Mayo 2026              \\[0.3cm]
        \textbf{Scripts:} & 6 archivos SQL         \\
    \end{tabular}
    \vfill
    {\small Universidad Nacional Aut\'onoma de M\'exico}
\end{titlepage}

\tableofcontents
\newpage

% ============================================================
\section{Descripci\'on General del Proyecto}

Este proyecto implementa el esquema de base de datos para el sistema de gesti\'on de un restaurante. Cubre la administraci\'on de empleados (cocineros, meseros y administrativos), cat\'alogo de productos (platillos y bebidas), clientes con facturaci\'on, \'ordenes y sus detalles.

El proyecto est\'a dividido en \textbf{seis archivos SQL}, cinco orientados a \textbf{PostgreSQL} y uno a \textbf{Oracle 23ai Free}:

\begin{center}
\begin{tabularx}{\textwidth}{llX}
    \toprule
    \textbf{Archivo} & \textbf{Motor} & \textbf{Prop\'osito} \\
    \midrule
    \texttt{00\_ejecutar\_todo.sql}        & Oracle 23ai & Script maestro completo \\
    \texttt{01\_ddl\_creacion\_tablas.sql} & PostgreSQL  & Creaci\'on de tablas \\
    \texttt{02\_secuencias.sql}            & PostgreSQL  & Definici\'on de secuencias \\
    \texttt{03\_triggers.sql}              & PostgreSQL  & L\'ogica de negocio con triggers \\
    \texttt{04\_funcion\_folio.sql}        & PostgreSQL  & Funci\'on generadora de folios \\
    \texttt{test\_triggers.sql}            & PostgreSQL  & Casos de prueba \\
    \bottomrule
\end{tabularx}
\end{center}

\subsection{Orden de ejecuci\'on en PostgreSQL}

\begin{enumerate}
    \item \texttt{02\_secuencias.sql}
    \item \texttt{01\_ddl\_creacion\_tablas.sql}
    \item \texttt{04\_funcion\_folio.sql}
    \item \texttt{03\_triggers.sql}
    \item \texttt{test\_triggers.sql}
\end{enumerate}

Las secuencias se ejecutan primero porque las tablas las referencian mediante \texttt{DEFAULT nextval(...)}.

% ============================================================
\section{Dise\~no del Esquema --- Estrategia Tabla por Subclase}

Se utiliza la estrategia \textit{tabla por subclase} (\textit{table per subclass}) para modelar la especializaci\'on de \texttt{EMPLEADO} y \texttt{PRODUCTO}.

\begin{itemize}
    \item Cada subtipo tiene su propia tabla con una FK que apunta a la tabla padre.
    \item Permite que un empleado tenga m\'ultiples roles simult\'aneamente (\textbf{especializaci\'on solapada}).
    \item Evita redundancia de datos y facilita la integridad referencial.
\end{itemize}

\subsection{Tablas del sistema}

\begin{center}
\begin{tabularx}{\textwidth}{llX}
    \toprule
    \textbf{Grupo} & \textbf{Tabla} & \textbf{Descripci\'on} \\
    \midrule
    Empleados     & \texttt{EMPLEADO}          & Entidad padre con datos generales \\
                  & \texttt{COCINERO}          & Subtipo: agrega especialidad \\
                  & \texttt{MESERO}            & Subtipo: agrega horario \\
                  & \texttt{ADMINISTRATIVO}    & Subtipo: agrega rol \\
    \midrule
    Extras        & \texttt{TELEFONO\_EMPLEADO} & Atributo multivaluado de empleado \\
                  & \texttt{DEPENDIENTE}        & Familiares del empleado \\
    \midrule
    Productos     & \texttt{PRODUCTO}           & Entidad padre con datos generales \\
                  & \texttt{PLATILLO}           & Subtipo de producto \\
                  & \texttt{BEBIDA}             & Subtipo de producto \\
    \midrule
    Cat\'alogo    & \texttt{CATEGORIA}          & Clasificaci\'on de productos \\
    \midrule
    \'Ordenes     & \texttt{ORDEN}              & Orden atendida por un mesero \\
                  & \texttt{DETALLE\_ORDEN}     & Renglones de la orden \\
    \midrule
    Facturaci\'on & \texttt{CLIENTE\_FACT}      & Clientes que solicitan factura \\
    \bottomrule
\end{tabularx}
\end{center}

% ============================================================
\section{02\_secuencias.sql --- Secuencias}

\subsection{Descripci\'on}

Define cuatro secuencias autoincrementales sin ciclo. Cada una genera llaves artificiales para una tabla espec\'ifica. Se ejecuta \textbf{antes} que el DDL de tablas porque estas las referencian.

\subsection{Explicaci\'on l\'inea por l\'inea}

\begin{lstlisting}[caption={Secuencia para folios de orden}]
DROP SEQUENCE IF EXISTS folio_seq;

CREATE SEQUENCE folio_seq
    START WITH 1      -- inicia en 1
    INCREMENT BY 1    -- sube de uno en uno
    MINVALUE 1        -- no puede bajar de 1
    NO MAXVALUE       -- sin limite superior
    NO CYCLE;         -- no reinicia al llegar al maximo
\end{lstlisting}

Genera los n\'umeros usados para construir folios tipo \texttt{ORD-001}.
El \texttt{NO CYCLE} es importante: garantiza que nunca se repita un folio aunque la secuencia llegue a valores muy altos.

\begin{lstlisting}[caption={Secuencias para llaves artificiales}]
DROP SEQUENCE IF EXISTS dependiente_id_seq;
CREATE SEQUENCE dependiente_id_seq
    START WITH 1 INCREMENT BY 1 MINVALUE 1 NO MAXVALUE NO CYCLE;

DROP SEQUENCE IF EXISTS categoria_id_seq;
CREATE SEQUENCE categoria_id_seq
    START WITH 1 INCREMENT BY 1 MINVALUE 1 NO MAXVALUE NO CYCLE;

DROP SEQUENCE IF EXISTS producto_id_seq;
CREATE SEQUENCE producto_id_seq
    START WITH 1 INCREMENT BY 1 MINVALUE 1 NO MAXVALUE NO CYCLE;
\end{lstlisting}

Las tres siguen el mismo patr\'on. Generan IDs para \texttt{DEPENDIENTE}, \texttt{CATEGORIA} y \texttt{PRODUCTO} respectivamente.

\begin{lstlisting}[caption={Verificaci\'on de secuencias creadas}]
SELECT sequencename, start_value, increment_by, min_value
FROM pg_sequences
WHERE schemaname = 'public'
ORDER BY sequencename;
\end{lstlisting}

Consulta al cat\'alogo interno de PostgreSQL (\texttt{pg\_sequences}) para confirmar que las cuatro secuencias quedaron correctamente creadas.

\subsection{Resumen de secuencias}

\begin{center}
\begin{tabular}{lll}
    \toprule
    \textbf{Secuencia} & \textbf{Usada en tabla} & \textbf{Prop\'osito} \\
    \midrule
    \texttt{folio\_seq}          & \texttt{ORDEN}       & N\'umero de folio de la orden \\
    \texttt{dependiente\_id\_seq} & \texttt{DEPENDIENTE} & PK de dependientes \\
    \texttt{categoria\_id\_seq}  & \texttt{CATEGORIA}   & PK de categor\'ias \\
    \texttt{producto\_id\_seq}   & \texttt{PRODUCTO}    & PK de productos \\
    \bottomrule
\end{tabular}
\end{center}

% ============================================================
\section{01\_ddl\_creacion\_tablas.sql --- Tablas}

\subsection{Descripci\'on}

Crea las 11 tablas del sistema en orden correcto (de independientes a dependientes, respetando las FK). Incluye una secci\'on de limpieza al inicio que elimina las tablas en orden inverso antes de crearlas.

\subsection{Limpieza previa}

\begin{lstlisting}[caption={DROP en orden inverso para respetar FKs}]
DROP TABLE IF EXISTS DETALLE_ORDEN    CASCADE;
DROP TABLE IF EXISTS ORDEN            CASCADE;
DROP TABLE IF EXISTS BEBIDA           CASCADE;
DROP TABLE IF EXISTS PLATILLO         CASCADE;
DROP TABLE IF EXISTS PRODUCTO         CASCADE;
DROP TABLE IF EXISTS CATEGORIA        CASCADE;
DROP TABLE IF EXISTS CLIENTE_FACT     CASCADE;
DROP TABLE IF EXISTS ADMINISTRATIVO   CASCADE;
DROP TABLE IF EXISTS COCINERO         CASCADE;
DROP TABLE IF EXISTS MESERO           CASCADE;
DROP TABLE IF EXISTS DEPENDIENTE      CASCADE;
DROP TABLE IF EXISTS TELEFONO_EMPLEADO CASCADE;
DROP TABLE IF EXISTS EMPLEADO         CASCADE;
\end{lstlisting}

El \texttt{IF EXISTS} evita error si la tabla no exist\'ia. El \texttt{CASCADE} elimina tambi\'en las restricciones que dependan de esa tabla.

\subsection{Tablas de empleados}

\begin{lstlisting}[caption={EMPLEADO --- entidad padre}]
CREATE TABLE EMPLEADO (
    num_empleado     VARCHAR(10)   NOT NULL,  -- PK alfanumerica (ej. EMP-001)
    nombre_pila      VARCHAR(40)   NOT NULL,
    apellido_paterno VARCHAR(40)   NOT NULL,
    apellido_materno VARCHAR(40)   NOT NULL,
    rfc              VARCHAR(13)   NOT NULL,  -- unico por empleado
    fecha_nacimiento DATE          NOT NULL,
    edad             INTEGER       NOT NULL,
    sueldo           NUMERIC(10,2) NOT NULL CHECK (sueldo > 0),
    foto             BYTEA,                  -- opcional, imagen binaria
    estado           VARCHAR(40)   NOT NULL,
    codigo_postal    VARCHAR(5)    NOT NULL,
    colonia          VARCHAR(100)  NOT NULL,
    calle            VARCHAR(100)  NOT NULL,
    numero           VARCHAR(10)   NOT NULL,
    CONSTRAINT pk_empleado     PRIMARY KEY (num_empleado),
    CONSTRAINT uq_empleado_rfc UNIQUE (rfc)
);
\end{lstlisting}

\begin{lstlisting}[caption={TELEFONO\_EMPLEADO --- atributo multivaluado}]
CREATE TABLE TELEFONO_EMPLEADO (
    num_empleado VARCHAR(10) NOT NULL,
    telefono     VARCHAR(15) NOT NULL,
    CONSTRAINT pk_telefono PRIMARY KEY (num_empleado, telefono),
    CONSTRAINT fk_telefono_empleado FOREIGN KEY (num_empleado)
        REFERENCES EMPLEADO (num_empleado)
        ON DELETE CASCADE ON UPDATE CASCADE
);
\end{lstlisting}

La PK compuesta \texttt{(num\_empleado, telefono)} permite que un empleado tenga varios tel\'efonos pero no el mismo n\'umero dos veces.

\begin{lstlisting}[caption={DEPENDIENTE --- familiares del empleado}]
CREATE TABLE DEPENDIENTE (
    dependiente_id INTEGER      NOT NULL DEFAULT nextval('dependiente_id_seq'),
    num_empleado   VARCHAR(10)  NOT NULL,
    nombre         VARCHAR(40)  NOT NULL,
    ap_paterno     VARCHAR(40)  NOT NULL,
    ap_materno     VARCHAR(40)  NOT NULL,
    curp           VARCHAR(18)  NOT NULL,  -- unico por dependiente
    parentesco     VARCHAR(30)  NOT NULL,
    CONSTRAINT pk_dependiente      PRIMARY KEY (dependiente_id),
    CONSTRAINT uq_dependiente_curp UNIQUE (curp),
    CONSTRAINT fk_dependiente_empleado FOREIGN KEY (num_empleado)
        REFERENCES EMPLEADO (num_empleado)
        ON DELETE CASCADE ON UPDATE CASCADE
);
\end{lstlisting}

\begin{lstlisting}[caption={Subtipos de EMPLEADO}]
CREATE TABLE COCINERO (
    num_empleado VARCHAR(10)  NOT NULL,
    especialidad VARCHAR(100) NOT NULL,
    CONSTRAINT pk_cocinero PRIMARY KEY (num_empleado),
    CONSTRAINT fk_cocinero_empleado FOREIGN KEY (num_empleado)
        REFERENCES EMPLEADO (num_empleado)
        ON DELETE CASCADE ON UPDATE CASCADE
);

CREATE TABLE MESERO (
    num_empleado VARCHAR(10) NOT NULL,
    horario      VARCHAR(50) NOT NULL,
    CONSTRAINT pk_mesero PRIMARY KEY (num_empleado),
    CONSTRAINT fk_mesero_empleado FOREIGN KEY (num_empleado)
        REFERENCES EMPLEADO (num_empleado)
        ON DELETE CASCADE ON UPDATE CASCADE
);

CREATE TABLE ADMINISTRATIVO (
    num_empleado VARCHAR(10) NOT NULL,
    rol          VARCHAR(50) NOT NULL,
    CONSTRAINT pk_administrativo PRIMARY KEY (num_empleado),
    CONSTRAINT fk_administrativo_empleado FOREIGN KEY (num_empleado)
        REFERENCES EMPLEADO (num_empleado)
        ON DELETE CASCADE ON UPDATE CASCADE
);
\end{lstlisting}

Cada subtipo solo agrega el atributo diferenciador. La FK con \texttt{ON DELETE CASCADE} garantiza que si se borra un empleado, sus registros de subtipo desaparecen tambi\'en.

\subsection{Tablas de productos}

\begin{lstlisting}[caption={CATEGORIA y PRODUCTO}]
CREATE TABLE CATEGORIA (
    categoria_id INTEGER NOT NULL DEFAULT nextval('categoria_id_seq'),
    nombre       VARCHAR(40) NOT NULL,
    descripcion  TEXT        NOT NULL,
    CONSTRAINT pk_categoria PRIMARY KEY (categoria_id)
);

CREATE TABLE PRODUCTO (
    producto_id    INTEGER       NOT NULL DEFAULT nextval('producto_id_seq'),
    nombre         VARCHAR(40)   NOT NULL,
    descripcion    TEXT          NOT NULL,
    precio         NUMERIC(8,2)  NOT NULL CHECK (precio > 0),
    disponibilidad BOOLEAN       NOT NULL DEFAULT TRUE,
    receta         TEXT          NOT NULL,
    categoria_id   INTEGER       NOT NULL,
    tipo           VARCHAR(10)   NOT NULL CHECK (tipo IN ('platillo', 'bebida')),
    CONSTRAINT pk_producto PRIMARY KEY (producto_id),
    CONSTRAINT fk_producto_categoria FOREIGN KEY (categoria_id)
        REFERENCES CATEGORIA (categoria_id)
        ON DELETE RESTRICT ON UPDATE CASCADE
);
\end{lstlisting}

El campo \texttt{tipo} es un discriminador que complementa a las tablas subtipo \texttt{PLATILLO} y \texttt{BEBIDA}. \texttt{disponibilidad} controla si el producto puede pedirse --- los triggers lo validan al insertar en \texttt{DETALLE\_ORDEN}.

\begin{lstlisting}[caption={Subtipos de PRODUCTO}]
CREATE TABLE PLATILLO (
    producto_id INTEGER NOT NULL,
    CONSTRAINT pk_platillo PRIMARY KEY (producto_id),
    CONSTRAINT fk_platillo_producto FOREIGN KEY (producto_id)
        REFERENCES PRODUCTO (producto_id)
        ON DELETE CASCADE ON UPDATE CASCADE
);

CREATE TABLE BEBIDA (
    producto_id INTEGER NOT NULL,
    CONSTRAINT pk_bebida PRIMARY KEY (producto_id),
    CONSTRAINT fk_bebida_producto FOREIGN KEY (producto_id)
        REFERENCES PRODUCTO (producto_id)
        ON DELETE CASCADE ON UPDATE CASCADE
);
\end{lstlisting}

\subsection{Tablas de \'ordenes}

\begin{lstlisting}[caption={CLIENTE\_FACT --- solo clientes con factura}]
CREATE TABLE CLIENTE_FACT (
    rfc              VARCHAR(13)  NOT NULL,  -- PK: identificador fiscal
    nombre_cliente   VARCHAR(40)  NOT NULL,
    apellido_paterno VARCHAR(40)  NOT NULL,
    apellido_materno VARCHAR(40)  NOT NULL,
    fecha_nacimiento DATE         NOT NULL,
    email            VARCHAR(60)  NOT NULL,
    estado           VARCHAR(40)  NOT NULL,
    codigo_postal    VARCHAR(5)   NOT NULL,
    colonia          VARCHAR(40)  NOT NULL,
    numero           VARCHAR(10)  NOT NULL,
    calle            VARCHAR(100) NOT NULL,
    razon_social     VARCHAR(150) NOT NULL,
    CONSTRAINT pk_cliente_fact PRIMARY KEY (rfc)
);
\end{lstlisting}

\begin{lstlisting}[caption={ORDEN --- una orden por mesero}]
CREATE TABLE ORDEN (
    folio        VARCHAR(10)   NOT NULL,  -- asignado por trigger
    fecha        DATE          NOT NULL DEFAULT CURRENT_DATE,
    hora         TIME          NOT NULL DEFAULT CURRENT_TIME,
    total        NUMERIC(10,2) NOT NULL DEFAULT 0.00,
    num_empleado VARCHAR(10)   NOT NULL,
    rfc          VARCHAR(13),             -- NULL si no requiere factura
    CONSTRAINT pk_orden PRIMARY KEY (folio),
    CONSTRAINT fk_orden_mesero FOREIGN KEY (num_empleado)
        REFERENCES MESERO (num_empleado)
        ON DELETE RESTRICT ON UPDATE CASCADE,
    CONSTRAINT fk_orden_cliente FOREIGN KEY (rfc)
        REFERENCES CLIENTE_FACT (rfc)
        ON DELETE SET NULL ON UPDATE CASCADE
);
\end{lstlisting}

El \texttt{ON DELETE SET NULL} en la FK del cliente permite que si se borra un cliente la orden siga existiendo, solo sin asociaci\'on de factura.

\begin{lstlisting}[caption={DETALLE\_ORDEN --- renglones de la orden}]
CREATE TABLE DETALLE_ORDEN (
    folio           VARCHAR(10)  NOT NULL,
    producto_id     INTEGER      NOT NULL,
    cantidad        INTEGER      NOT NULL CHECK (cantidad > 0),
    precio_platillo NUMERIC(8,2) NOT NULL CHECK (precio_platillo > 0),
    CONSTRAINT pk_detalle_orden    PRIMARY KEY (folio, producto_id),
    CONSTRAINT fk_detalle_orden    FOREIGN KEY (folio)
        REFERENCES ORDEN (folio)
        ON DELETE CASCADE ON UPDATE CASCADE,
    CONSTRAINT fk_detalle_producto FOREIGN KEY (producto_id)
        REFERENCES PRODUCTO (producto_id)
        ON DELETE RESTRICT ON UPDATE CASCADE
);
\end{lstlisting}

La PK compuesta \texttt{(folio, producto\_id)} garantiza que un producto no se repita dentro de la misma orden. El \texttt{precio\_platillo} es calculado autom\'aticamente por el trigger --- el usuario puede enviar \texttt{0}.

% ============================================================
\section{04\_funcion\_folio.sql --- Funci\'on generar\_folio}

\subsection{Descripci\'on}

Funci\'on auxiliar que consume \texttt{folio\_seq} y devuelve un folio formateado. Es invocada por el Trigger 1 cada vez que se inserta una orden.

\subsection{Explicaci\'on l\'inea por l\'inea}

\begin{lstlisting}[caption={Funci\'on generar\_folio()}]
CREATE OR REPLACE FUNCTION generar_folio()
RETURNS VARCHAR(10) AS $$
-- No recibe parametros. Retorna VARCHAR(10),
-- suficiente para ORD-001 hasta ORD-9999.
DECLARE
    v_numero INTEGER;     -- guarda el numero de la secuencia
    v_folio  VARCHAR(10); -- guarda el folio formateado
BEGIN
    -- Obtiene el siguiente valor de la secuencia.
    -- Cada llamada incrementa el contador y nunca repite.
    v_numero := nextval('folio_seq');

    -- Construye el folio en tres pasos:
    -- 1. v_numero::TEXT  -> convierte entero a texto  (1 -> '1')
    -- 2. LPAD(..., 3, '0') -> rellena con ceros       ('1' -> '001')
    -- 3. 'ORD-' ||       -> concatena el prefijo  ('ORD-001')
    -- Si el numero supera 3 digitos, LPAD no recorta:
    -- 1000 -> '1000' -> 'ORD-1000' (7 caracteres, entra en VARCHAR(10))
    v_folio := 'ORD-' || LPAD(v_numero::TEXT, 3, '0');

    RETURN v_folio;
END;
$$ LANGUAGE plpgsql;
\end{lstlisting}

\subsection{Ejemplos de salida}

\begin{center}
\begin{tabular}{cc}
    \toprule
    \textbf{Valor de secuencia} & \textbf{Folio generado} \\
    \midrule
    1    & \texttt{ORD-001}  \\
    50   & \texttt{ORD-050}  \\
    99   & \texttt{ORD-099}  \\
    100  & \texttt{ORD-100}  \\
    1000 & \texttt{ORD-1000} \\
    \bottomrule
\end{tabular}
\end{center}

% ============================================================
\section{03\_triggers.sql --- Triggers}

\subsection{Descripci\'on general}

Cuatro triggers que automatizan la l\'ogica de negocio sobre \texttt{ORDEN} y \texttt{DETALLE\_ORDEN}. En PostgreSQL cada trigger requiere una funci\'on separada.

\begin{center}
\begin{tabularx}{\textwidth}{llllX}
    \toprule
    \textbf{\#} & \textbf{Trigger} & \textbf{Tabla} & \textbf{Momento} & \textbf{Acci\'on} \\
    \midrule
    1 & \texttt{trg\_generar\_folio}     & \texttt{ORDEN}        & BEFORE INSERT & Asigna folio autom\'atico \\
    2 & \texttt{trg\_insertar\_detalle}  & \texttt{DETALLE\_ORDEN} & BEFORE INSERT & Valida, calcula subtotal, actualiza total \\
    3 & \texttt{trg\_actualizar\_detalle}& \texttt{DETALLE\_ORDEN} & BEFORE UPDATE & Recalcula subtotal y total \\
    4 & \texttt{trg\_eliminar\_detalle}  & \texttt{DETALLE\_ORDEN} & AFTER DELETE  & Descuenta subtotal del total \\
    \bottomrule
\end{tabularx}
\end{center}

\subsection{Trigger 1 --- Folio autom\'atico}

\begin{lstlisting}[caption={Funci\'on y trigger de folio autom\'atico}]
CREATE OR REPLACE FUNCTION fn_generar_folio_orden()
RETURNS TRIGGER AS $$
BEGIN
    -- Asigna el folio generado al registro nuevo antes de insertarlo.
    -- NEW contiene los valores que se van a insertar.
    NEW.folio := generar_folio();
    RETURN NEW;  -- obligatorio en triggers BEFORE: devuelve el registro modificado
END;
$$ LANGUAGE plpgsql;

DROP TRIGGER IF EXISTS trg_generar_folio ON ORDEN;

CREATE TRIGGER trg_generar_folio
    BEFORE INSERT ON ORDEN  -- se ejecuta antes del INSERT
    FOR EACH ROW            -- una vez por cada fila insertada
    EXECUTE FUNCTION fn_generar_folio_orden();
\end{lstlisting}

El usuario no debe proporcionar el folio. Si lo proporciona, el trigger lo sobreescribe.

\subsection{Trigger 2 --- Insertar en DETALLE\_ORDEN}

\begin{lstlisting}[caption={Funci\'on del trigger de inserci\'on}]
CREATE OR REPLACE FUNCTION fn_insertar_detalle_orden()
RETURNS TRIGGER AS $$
DECLARE
    v_disponible BOOLEAN;
    v_precio     NUMERIC(8,2);
    v_nombre     VARCHAR(40);
BEGIN
    -- (a) Buscar el producto en catalogo
    SELECT disponibilidad, precio, nombre
    INTO v_disponible, v_precio, v_nombre
    FROM PRODUCTO
    WHERE producto_id = NEW.producto_id;

    -- Si el SELECT no encontro nada, FOUND = false
    IF NOT FOUND THEN
        RAISE EXCEPTION 'El producto con ID % no existe.', NEW.producto_id;
    END IF;

    -- Validar que este disponible para venta
    IF v_disponible = FALSE THEN
        RAISE EXCEPTION 'El producto "%" no esta disponible actualmente.', v_nombre;
    END IF;

    -- (b) Calcular subtotal: cantidad x precio unitario
    -- El usuario puede mandar precio_platillo = 0; el trigger lo corrige
    NEW.precio_platillo := NEW.cantidad * v_precio;

    -- (c) Actualizar total de la orden
    -- COALESCE evita NULL cuando la orden aun no tiene renglones
    UPDATE ORDEN
    SET total = (
        SELECT COALESCE(SUM(precio_platillo), 0) + NEW.precio_platillo
        FROM DETALLE_ORDEN
        WHERE folio = NEW.folio
    )
    WHERE folio = NEW.folio;

    RETURN NEW;
END;
$$ LANGUAGE plpgsql;
\end{lstlisting}

\subsection{Trigger 3 --- Actualizar en DETALLE\_ORDEN}

\begin{lstlisting}[caption={Funci\'on del trigger de actualizaci\'on}]
CREATE OR REPLACE FUNCTION fn_actualizar_detalle_orden()
RETURNS TRIGGER AS $$
DECLARE
    v_precio NUMERIC(8,2);
BEGIN
    -- Obtener precio actual (puede haber cambiado desde el INSERT)
    SELECT precio INTO v_precio
    FROM PRODUCTO
    WHERE producto_id = NEW.producto_id;

    -- Recalcular subtotal con la nueva cantidad
    NEW.precio_platillo := NEW.cantidad * v_precio;

    -- Recalcular total de la orden.
    -- El CASE es necesario: el UPDATE en DETALLE_ORDEN aun no se
    -- aplico cuando el trigger corre, entonces el renglon modificado
    -- todavia tiene el valor viejo en la tabla. Usamos NEW.precio_platillo
    -- para ese renglon y el valor almacenado para los demas.
    UPDATE ORDEN
    SET total = (
        SELECT COALESCE(SUM(
            CASE
                WHEN folio = NEW.folio AND producto_id = NEW.producto_id
                THEN NEW.precio_platillo   -- valor nuevo calculado
                ELSE precio_platillo       -- valor existente en tabla
            END
        ), 0)
        FROM DETALLE_ORDEN
        WHERE folio = NEW.folio
    )
    WHERE folio = NEW.folio;

    RETURN NEW;
END;
$$ LANGUAGE plpgsql;
\end{lstlisting}

\subsection{Trigger 4 --- Eliminar de DETALLE\_ORDEN}

\begin{lstlisting}[caption={Funci\'on del trigger de eliminaci\'on}]
CREATE OR REPLACE FUNCTION fn_eliminar_detalle_orden()
RETURNS TRIGGER AS $$
BEGIN
    -- Recalcular total excluyendo el renglon eliminado.
    -- Este trigger es AFTER DELETE, entonces el renglon ya no existe
    -- en la tabla cuando se ejecuta esta funcion.
    -- OLD contiene los valores del renglon que se elimino.
    UPDATE ORDEN
    SET total = (
        SELECT COALESCE(SUM(precio_platillo), 0)
        FROM DETALLE_ORDEN
        WHERE folio = OLD.folio
          AND producto_id != OLD.producto_id  -- excluir el eliminado
    )
    WHERE folio = OLD.folio;

    RETURN OLD;  -- en triggers AFTER DELETE se retorna OLD
END;
$$ LANGUAGE plpgsql;

CREATE TRIGGER trg_eliminar_detalle
    AFTER DELETE ON DETALLE_ORDEN  -- AFTER porque necesitamos que ya este borrado
    FOR EACH ROW
    EXECUTE FUNCTION fn_eliminar_detalle_orden();
\end{lstlisting}

\textbf{Nota importante:} Este trigger usa \texttt{AFTER DELETE} a diferencia de los anteriores que usan \texttt{BEFORE}. Esto es necesario porque necesita que el registro ya est\'e eliminado de la tabla para que el \texttt{SUM} lo excluya de forma natural.

% ============================================================
\section{test\_triggers.sql --- Casos de Prueba}

\subsection{Descripci\'on}

Siete casos de prueba que verifican el correcto funcionamiento de los cuatro triggers. Incluye datos de prueba, operaciones y resultados esperados.

\subsection{Datos base}

\begin{lstlisting}[caption={Inserci\'on de datos de prueba}]
-- Categoria
INSERT INTO CATEGORIA (categoria_id, nombre, descripcion)
VALUES (nextval('categoria_id_seq'), 'Comida mexicana', 'Platillos tipicos mexicanos');

-- Empleado
INSERT INTO EMPLEADO (num_empleado, nombre_pila, apellido_paterno, apellido_materno,
    rfc, fecha_nacimiento, edad, sueldo, estado, codigo_postal, colonia, calle, numero)
VALUES ('EMP-001', 'Juan', 'Garcia', 'Lopez',
    'GALJ900101ABC', '1990-01-01', 34, 8000.00,
    'CDMX', '06600', 'Juarez', 'Reforma', '100');

-- Mesero (subtipo de empleado)
INSERT INTO MESERO (num_empleado, horario)
VALUES ('EMP-001', 'Lunes a Viernes 08:00-16:00');

-- Productos: 2 disponibles, 1 no disponible
INSERT INTO PRODUCTO (producto_id, nombre, descripcion, precio,
                      disponibilidad, receta, categoria_id, tipo)
VALUES
    (nextval('producto_id_seq'), 'Tacos de pastor',
     'Tacos con carne al pastor',   85.00, TRUE,
     'Carne de cerdo marinada...', 1, 'platillo'),
    (nextval('producto_id_seq'), 'Agua de jamaica',
     'Bebida de flor de jamaica',   25.00, TRUE,
     'Flor de jamaica, agua, azucar...', 1, 'bebida'),
    (nextval('producto_id_seq'), 'Enchiladas verdes',
     'Enchiladas con salsa verde',  95.00, FALSE,  -- no disponible
     'Tortillas, salsa verde...', 1, 'platillo');

INSERT INTO PLATILLO (producto_id) VALUES (1);
INSERT INTO BEBIDA   (producto_id) VALUES (2);
INSERT INTO PLATILLO (producto_id) VALUES (3);
\end{lstlisting}

\subsection{Pruebas}

\begin{lstlisting}[caption={Prueba 1: Folio autom\'atico}]
INSERT INTO ORDEN (fecha, hora, num_empleado)
VALUES (CURRENT_DATE, CURRENT_TIME, 'EMP-001');
-- El usuario NO proporciona folio; el trigger lo asigna.

SELECT folio, fecha, hora, total FROM ORDEN;
-- Resultado esperado: folio = 'ORD-001', total = 0.00
\end{lstlisting}

\begin{lstlisting}[caption={Prueba 2: INSERT con producto disponible}]
-- 2 tacos de pastor a $85 c/u = $170
-- Se envia precio_platillo = 0; el trigger lo corrige
INSERT INTO DETALLE_ORDEN (folio, producto_id, cantidad, precio_platillo)
VALUES ('ORD-001', 1, 2, 0);

SELECT * FROM DETALLE_ORDEN WHERE folio = 'ORD-001';
-- Resultado esperado: precio_platillo = 170.00

SELECT folio, total FROM ORDEN WHERE folio = 'ORD-001';
-- Resultado esperado: total = 170.00
\end{lstlisting}

\begin{lstlisting}[caption={Prueba 3: Acumulaci\'on de total}]
-- 1 agua de jamaica a $25
INSERT INTO DETALLE_ORDEN (folio, producto_id, cantidad, precio_platillo)
VALUES ('ORD-001', 2, 1, 0);

SELECT folio, total FROM ORDEN WHERE folio = 'ORD-001';
-- Resultado esperado: total = 195.00  ($170 + $25)
\end{lstlisting}

\begin{lstlisting}[caption={Prueba 4: UPDATE de cantidad}]
-- Cambiar tacos de 2 a 3 unidades: 3 x $85 = $255
UPDATE DETALLE_ORDEN SET cantidad = 3
WHERE folio = 'ORD-001' AND producto_id = 1;

SELECT * FROM DETALLE_ORDEN WHERE folio = 'ORD-001';
-- Resultado esperado: precio_platillo de tacos = 255.00

SELECT folio, total FROM ORDEN WHERE folio = 'ORD-001';
-- Resultado esperado: total = 280.00  ($255 + $25)
\end{lstlisting}

\begin{lstlisting}[caption={Prueba 5: DELETE de un rengl\'on}]
-- Eliminar agua de jamaica
DELETE FROM DETALLE_ORDEN
WHERE folio = 'ORD-001' AND producto_id = 2;

SELECT folio, total FROM ORDEN WHERE folio = 'ORD-001';
-- Resultado esperado: total = 255.00  (se restan los $25)
\end{lstlisting}

\begin{lstlisting}[caption={Prueba 6: Producto no disponible}]
-- Enchiladas verdes tiene disponibilidad = FALSE
-- El bloque DO captura la excepcion sin abortar el script
DO $$
BEGIN
    INSERT INTO DETALLE_ORDEN (folio, producto_id, cantidad, precio_platillo)
    VALUES ('ORD-001', 3, 1, 0);

    -- Esta linea solo se ejecuta si el trigger falla en lanzar el error
    RAISE NOTICE 'ERROR: El trigger no lanzo la excepcion esperada';
EXCEPTION
    WHEN OTHERS THEN
        -- SQLERRM contiene el mensaje de error del trigger
        RAISE NOTICE 'OK: Excepcion capturada correctamente -> %', SQLERRM;
END;
$$;
-- Resultado esperado: NOTICE con el mensaje del trigger
\end{lstlisting}

\begin{lstlisting}[caption={Prueba 7: Segunda orden con folio incremental}]
INSERT INTO ORDEN (fecha, hora, num_empleado)
VALUES (CURRENT_DATE, CURRENT_TIME, 'EMP-001');

SELECT folio FROM ORDEN ORDER BY folio;
-- Resultado esperado: ORD-001 y ORD-002
\end{lstlisting}

\subsection{Resumen de pruebas}

\begin{center}
\begin{tabularx}{\textwidth}{clX}
    \toprule
    \textbf{Prueba} & \textbf{Operaci\'on} & \textbf{Resultado esperado} \\
    \midrule
    1 & INSERT en ORDEN         & folio = \texttt{ORD-001}, total = \$0.00 \\
    2 & INSERT 2 tacos (\$85)   & subtotal = \$170, total = \$170 \\
    3 & INSERT 1 jamaica (\$25) & total = \$195 \\
    4 & UPDATE tacos a 3 uds.   & subtotal = \$255, total = \$280 \\
    5 & DELETE jamaica          & total = \$255 \\
    6 & INSERT no disponible    & excepci\'on controlada \\
    7 & Segunda orden           & folios \texttt{ORD-001} y \texttt{ORD-002} \\
    \bottomrule
\end{tabularx}
\end{center}

% ============================================================
\section{Diferencias entre la versi\'on Oracle y PostgreSQL}

\begin{center}
\begin{tabularx}{\textwidth}{lXX}
    \toprule
    \textbf{Aspecto} & \textbf{Oracle 23ai (00\_ejecutar\_todo)} & \textbf{PostgreSQL (01--04)} \\
    \midrule
    Tipo binario     & \texttt{BLOB}            & \texttt{BYTEA} \\
    Variables en trigger & \texttt{:NEW.campo}, \texttt{:OLD.campo} & \texttt{NEW.campo}, \texttt{OLD.campo} \\
    ON UPDATE CASCADE & No soportado (omitido)  & S\'i soportado \\
    Estructura trigger & C\'odigo directo en el trigger & Requiere funci\'on separada (\texttt{fn\_...}) \\
    Problema tabla mutante & Soluci\'on incremental (\texttt{total = total + subtotal}) & No aplica; se usa \texttt{SUM()} completo \\
    Tipo hora        & \texttt{DATE} (incluye hora) con \texttt{DEFAULT SYSDATE} & \texttt{TIME} separado con \texttt{DEFAULT CURRENT\_TIME} \\
    Manejo de errores & \texttt{RAISE\_APPLICATION\_ERROR(-20001, msg)} & \texttt{RAISE EXCEPTION 'msg'} \\
    Verificar errores & \texttt{SHOW ERRORS}     & \texttt{NOT FOUND} (variable impl\'icita) \\
    \bottomrule
\end{tabularx}
\end{center}

\end{document}
